\section{Related Work}
\label{sec:related}

\subsection{UAV Voice Interfaces Under Rotor Noise}

Prior work has studied speech interfaces for unmanned aerial vehicles, showing that voice can be a feasible interaction channel but also that the UAV platform creates a difficult acoustic setting. Kite explores automatic speech recognition for UAV control \cite{oneata19_interspeech}, and subsequent multimodal work extends UAV speech recognition beyond a single audio stream \cite{oneata2021multimodal}. These systems motivate the premise that voice can be useful around drones. They also leave open a systems question that is central to this paper: how the output of a noisy voice interface should be constrained before it can affect vehicle behavior.

Other UAV-focused studies show why this question is difficult. Drone operating noise can degrade voice-command discrimination \cite{miesikowska2021uavnoise}, and speech recognition for UAV control remains sensitive to the acoustic and operating context \cite{park2023uavvehicle}. Drone audition and sound-source-localization surveys similarly identify ego-noise, payload limits, and onboard real-time processing as recurring obstacles for aerial audio sensing \cite{chevtchenko2025drone_ssl}. Our work takes this platform-generated noise condition as the default operating environment rather than a peripheral robustness test. The goal is not only to improve recognition under noise, but to use recognition as one component of an additional UAV safety layer whose output remains mediated by a conservative control boundary.

\subsection{From Keywords to Spoken Intent Under Constraints}

Small-footprint keyword spotting is the closest mature line of on-device speech work. Compact CNNs, microcontroller-oriented keyword models, public limited-vocabulary datasets, and streaming keyword systems have shown that always-on speech interfaces must be evaluated with false accepts, misses, latency, and memory rather than accuracy alone \cite{sainath2015cnn_kws,warden2018speechcommands,zhang2017helloedge,rybakov2020streamingkws}. This literature is useful because it treats speech recognition as an online decision process. It is still not the same problem as a UAV safety interaction layer: a detected keyword can be a trigger, but a nearby drone needs a mediated state event before any physical action is considered.

Speech-to-intent models provide an alternative to transcript-first pipelines when the system needs a compact action abstraction rather than a full word sequence. End-to-end spoken language understanding has shown that useful intent decisions can be learned directly from speech \cite{lugosch19_interspeech,kuo20_interspeech,mhiri20_interspeech}, including streaming and on-device-motivated settings \cite{ijcai2021p538,le2022deliberation}. General spoken language understanding resources such as SLURP and SLUE broaden the available evaluation space for speech intent tasks \cite{bastianelli-etal-2020-slurp,shon2022slue}, and multilingual benchmarks such as ML-SUPERB motivate speaker and language robustness questions \cite{shi23g_interspeech}. These benchmarks are useful context, but they do not directly measure rotor self-noise, embedded deployment, or safety-state behavior. In this paper, speaker and language variation are therefore treated as evaluation axes rather than completed generalization claims.

\subsection{Robust Embedded Audio and Safety-Oriented Evaluation}

Robust speech learning and frontend techniques are relevant because the target condition contains severe self-noise. Multi-style training and gain control improve small-footprint keyword spotting under noisy and far-field conditions \cite{prabhavalkar2015agc}; PCEN moves part of loudness and far-field robustness into a trainable frontend \cite{wang2017pcen}; and SpecAugment improves robustness through time-frequency masking without changing the inference model \cite{park2019specaugment}. Noisy-student training and self-supervised speech representations have improved robustness in broader speech settings \cite{park2020improvednoisystudent,baevski2020wav2vec20,yang21c_interspeech}. We build on this family of ideas through a clean-to-noisy training recipe, but keep the claim local: the recognizer is evaluated as a rotor-noisy UAV intent recognizer, not as a general-purpose robust speech model.

Embedded speech systems are often evaluated as standalone keyword or command recognizers, but deployment adds constraints that are not visible from model accuracy alone. TinyML runtimes and kernels such as TensorFlow Lite Micro and CMSIS-NN make microcontroller inference practical, while MCUNet and MLPerf Tiny show that memory, operator support, latency, energy, and real-hardware measurements are first-order design variables \cite{david2021tflm,lai2018cmsisnn,lin2020mcunet,banbury2021mlperftiny}. A UAV safety interaction layer requires an even broader evidence structure: the board must report timing and confidence, the bridge must log state transitions, and movement-like outputs must remain non-command events unless a downstream policy permits action. Our deployment contribution is therefore not simply an embedded model artifact. It is the path from onboard audio capture to an auditable voice event that can be evaluated separately for recognition quality, runtime feasibility, response time, and safety-state behavior.
